% =====================================================================
%  Chapter 12 — The alignment pincer (S32–S38)
% =====================================================================
\chapter{The Alignment Pincer}
\label{ch:pincer}

\begin{record}
\textbf{What this chapter covers.} Sessions \Sn{32}--\Sn{38}: the programme
rebuilt after the audit. Its object is a scalar, \(\sstar\), measuring how
disordered the vorticity direction field is at the self-similar scale. Its
strategy is a pincer: two named hypotheses, H1 and H2, which together forbid a
Type-I singularity from ever becoming coherent. Its result, at the end of the
arc, is that the one alternative escape route is \IMPOSSIBLE{}, and therefore
that the hypothesis the pincer needs is not one option among several but
\emph{necessary}. That is a sharpening of the problem, and it is also a
narrowing of the programme's own options. Both readings are correct and both
are recorded.

\smallskip
Sessions \Sn{36} and \Sn{37} belong to this arc but are treated separately, in
Chapter~\ref{ch:unification}, because their content --- verification of a
derived identity, and the collapse of two open problems into one --- is of a
different kind from the construction described here.
\end{record}

\section{The starting position}
\label{sec:pincer-start}

Chapter~\ref{ch:inheritance} established what \Sn{32} had to work with. The
\Sn{31} audit left two sound objects on the geometric side: the orthogonal
split identity \(\abs{\nabla\omega}^2=\abs{\nabla m}^2+m^2\abs{\nabla\ehat}^2\)
(\texttt{lem:grad-omega-split}, line~4564) and the exact scale invariance of
\(\sigma=A_{\mathrm{loc}}/M^{3/2}\). It left one defect that later became the
programme's central obstruction. Everything else on that side --- the
coercivity, the Gr\"onwall closure, Theorem~E --- was withdrawn.

\Sn{32} rebuilt on the two survivors and nothing else. That constraint is the
reason the chapter's first three constructions are \emph{re-derivations of
things the withdrawn work had asserted}, and it is the reason they are sound.

\subsection{The direction equation, derived rather than assumed}

Write \(m=\abs{\omega}\), \(\ehat=\omega/\abs{\omega}\),
\(\alpha=\ehat\cdot S\ehat\) for the stretching rate, \(\Dt=\partial_t+u\cdot\nabla\),
and \(w=\abs{\nabla\ehat}^2\). Constantin's pair of equations
(\texttt{lem:direction-eq}, line~5349) is:
\begin{align}
  \Dt m &= \alpha m + \nu\bigl(\Delta m - m\,w\bigr),
    \label{eq:pincer-mag}\\
  \Dt\ehat &= \nu\Delta\ehat + 2\nu(\nabla\log m\cdot\nabla)\ehat
    + \nu\,w\,\ehat + P_{\ehat^{\perp}}\bigl((\ehat\cdot\nabla)u\bigr).
    \label{eq:pincer-dir}
\end{align}

Two features of this pair do all the work in the sessions that follow. First,
\emph{no pressure appears}. The document says so in terms: this is the same
structural advantage the withdrawn \S20 was reaching for, obtained without any
invalid estimate. Second, the magnitude equation \eqref{eq:pincer-mag} carries
a good-sign sink \(-\nu m w\) --- the radial equation dissipates exactly the
quantity the angular equation needs controlled. \Sn{34} and \Sn{35} are both
built on that observation.

\subsection{\texorpdfstring{The coherence deficit \(\sstar\)}{The coherence deficit sigma-star}}

\begin{defn}[\texttt{def:sigma-star}, line~5393]
\label{def:sstar}
Let \(M(t)=\norm{\omega(\cdot,t)}_{L^\infty}\), let \(\rstar=M^{-1/2}\) be the
self-similar scale, and let \(x^*\) be a point with \(\abs{\omega(x^*)}\ge
M/2\). Then
\[
  \sstar(t) \;=\; (\rstar)^{-1}\int_{B_{\rstar}(x^*)}
    \abs{\nabla\ehat}^2\,\mathbf 1_{\{\abs{\omega}\ge M/4\}}\dee x .
\]
\end{defn}

\texttt{prop:sigma-star-invariant} (line~5402) proves that this is exactly
invariant under the Navier--Stokes scaling group, and the calibration remark
that follows fixes what its size means: if \(\abs{\nabla\ehat}\sim1/\rstar\)
throughout the ball --- borderline Constantin--Fefferman coherence at the
self-similar scale --- then \(\sstar\sim1\). So \(\sstar\ll1\) is
\emph{better} than borderline, and \(\sstar\) is the natural small parameter
for an \(\varepsilon\)-regularity statement about the direction field.

\texttt{prop:sigma-controls} (line~5424) proves \(\sstar\le16\sigma\) in four
lines. That inequality is the formal moment of inheritance: the quantity the
audit certified controls the quantity the new programme uses.

\subsection{\texorpdfstring{Why \(\rstar\) is the right scale: marginality}{Why r-star is the right scale: marginality}}

\texttt{rem:marginality} (line~5443) is short and is the structural reason the
route was thought plausible at all. In any parabolic estimate for
\eqref{eq:pincer-dir} on the cylinder
\(Q(\rstar)=B_{\rstar}\times[t_0-(\rstar)^2,t_0]\), the stretching term acts as
a potential of size \(\norm{\nabla u}_{L^\infty}\). Over the time span
\((\rstar)^2=M^{-1}\), its Gr\"onwall exponent is
\[
  \norm{\nabla u}_{L^\infty}\cdot(\rstar)^2
  \;\lesssim\; M\cdot M^{-1}\cdot\Lfac \;=\; \Lfac .
\]
The direction equation is therefore \emph{exactly critical} at scale
\(\rstar\), up to a logarithm --- where the corresponding statement for \(u\)
itself, the object CKN treats, is supercritical. This is the whole appeal of
working with \(\ehat\) rather than \(u\).

It is also where the logarithm enters, and \texttt{rem:marginality} says so
explicitly, naming its source: ``the logarithm is mandatory:
Calder\'on--Zygmund operators are unbounded on \(L^\infty\) --- this was
Defect~8 of the S31 audit''. Chapter~\ref{ch:cz-log} follows that factor
through the whole programme. For this chapter it is enough to note that the
marginality which makes the route plausible and the logarithm which obstructs
it are \emph{the same fact}, read twice.

\section{The pincer itself}
\label{sec:pincer-statement}

\begin{statement}[H1, the Bridge Lemma --- \texttt{conj:bridge}, line~5466]
\label{st:h1}
There exist \(\varepsilon_0>0\) and \(C_B\ge1\) (possibly degrading
logarithmically in \(M\)) such that, if \(u\) satisfies the Type-I bound and
\(\sup_{t}\sstar(t)\le\varepsilon_0\) over the cylinder, then
\[
  \norm{\nabla\ehat}_{L^\infty\left(Q(\rstar/2)\cap\{\abs{\omega}\ge M/2\}\right)}
  \;\le\; \frac{C_B}{\rstar}.
\]
Status: \OPEN{} target. It upgrades \emph{averaged} coherence to
\emph{pointwise} coherence.
\end{statement}

\begin{statement}[H2, localized Constantin--Fefferman]
\label{st:h2}
Pointwise coherence \(\norm{\nabla\ehat}_{L^\infty}\le C_B/\rstar\) on
\(Q(\rstar/2)\cap\{\abs{\omega}\ge M/2\}\) implies the depletion bound
\(\alpha(x^*,t)\le C M^{1/2+\theta}\) for some \(\theta<1/2\) near a putative
Type-I blow-up time. Status: \OPEN{} target at \Sn{32}; derived in \Sn{37} and
found to split into two regimes (Chapter~\ref{ch:unification}).
\end{statement}

\begin{statement}[The pincer --- \texttt{thm:pincer-conditional}, line~5495]
\label{st:pincer}
\CONDITIONAL{H1, H2} No Type-I singularity can satisfy
\(\liminf_{t\to T^-}\sstar(t)\le\varepsilon_0\). Every Type-I singularity must
maintain \emph{persistent geometric disorder}, \(\sstar(t)>\varepsilon_0\), for
all \(t\) close to \(T\).
\end{statement}

The proof given H1 and H2 is three lines and is the part of the construction
nobody has ever disputed. If \(\sstar(t_k)\le\varepsilon_0\) along
\(t_k\to T^-\), then H1 gives pointwise coherence, H2 converts it into
\(\tfrac{\dee}{\dee t}\log M=\alpha(x^*,t)\le CM^{1/2+\theta}\) with
\(1/2+\theta<1\), and \(\int^T M^{1/2+\theta}\dee t\) converges under
\(M\le C_I/(T-t)\) precisely when \(1/2+\theta<1\). So \(M\) stays bounded,
contradicting blow-up.

\begin{figure}[ht]
\centering
\begin{tikzpicture}[
  >={Latex[length=2mm]},
  bx/.style={draw, rounded corners=2pt, font=\scriptsize, align=center,
             inner sep=3.5pt, text width=33mm, minimum height=10mm},
  hyp/.style={bx, draw=openblue, thick, fill=panelgrey},
  prv/.style={bx, draw=provedgreen, very thick},
  cnc/.style={bx, draw=impossiblered, thick},
  ar/.style={->, thick}
]
\node[bx] (assume) at (0,0)
  {\textbf{Assume:} a Type-I singularity at \(T\), and\\
   \(\liminf_{t\to T^-}\sstar(t)\le\varepsilon_0\)};

\node[hyp] (h1) at (4.9,1.5)
  {\textbf{H1} Bridge Lemma\\ averaged \(\Rightarrow\) pointwise\\
   \(\norm{\nabla\ehat}_\infty\le C_B/\rstar\)};
\node[hyp] (h2) at (4.9,-1.5)
  {\textbf{H2} localized CF\\ coherence \(\Rightarrow\) depletion\\
   \(\alpha\le CM^{1/2+\theta}\), \(\theta<\tfrac12\)};

\node[bx] (int) at (9.6,0)
  {\(\dfrac{\dee}{\dee t}\log M\le CM^{1/2+\theta}\)\\[2pt]
   integrable under Type~I};
\node[cnc] (contra) at (9.6,-2.7)
  {\(M\) stays bounded:\\ \textbf{contradiction}};

\draw[ar] (assume) -- (h1);
\draw[ar] (h1) -- (h2);
\draw[ar] (h2) -- (int);
\draw[ar] (int) -- (contra);

\node[prv, below=17mm of assume] (s38)
  {\Sn{38}: Regime~A \IMPOSSIBLE\\
   \(\Rightarrow\) \(\sstar\)-decay is \emph{necessary}};
\draw[ar, provedgreen, dashed] (s38.north) -- (assume.south);

\node[font=\scriptsize\itshape, text=rulegrey, align=left, text width=52mm]
  at (4.6,-4.0)
  {The conclusion is \emph{persistent disorder}, not regularity. Closing the
   pincer excludes coherent Type-I singularities only; the remaining
   possibility --- a singularity that stays disordered --- is what
   \(\sstar\)-decay would have to rule out.};
\end{tikzpicture}
\caption[The alignment pincer]{The pincer of \texttt{thm:pincer-conditional}
(line~5495). Two hypotheses close on the assumption of eventual coherence and
produce a contradiction with blow-up. \Sn{38}'s contribution, shown dashed, is
that the disjunction the argument sits inside has only one other branch, and
that branch is \IMPOSSIBLE{} --- which makes \(\sstar\)-decay necessary rather
than convenient.}
\label{fig:pincer}
\end{figure}

\section{H1 attacked: the level-set Moser scheme}
\label{sec:moser}

H1 is a Moser iteration for \(w=\abs{\nabla\ehat}^2\), which by
\eqref{eq:pincer-dir} satisfies a drift-diffusion inequality. Two obstacles
were named at \Sn{32} and neither is cosmetic.

\begin{enumerate}[leftmargin=2.2em]
\item \emph{The singular drift.} The drift is \(u+2\nu\nabla\log m\). On
  \(\{\abs{\omega}\ge M/4\}\) one has \(\abs{\nabla\log m}\le4\abs{\nabla\omega}/M\),
  which requires a pointwise gradient bound that is not available.
\item \emph{The marginal potential.} By \texttt{rem:marginality} the potential
  is critical up to the logarithm, so the iteration constants degrade
  logarithmically.
\end{enumerate}

\subsection{\texorpdfstring{\Sn{33}: the first rung}{S33: the first rung}}

The device that disposes of obstacle~(i) is to run the iteration on
\emph{level sets} of \(\abs{\omega}\), where the drift enters only through its
divergence structure. \Sn{33} builds a composite cutoff
\(\eta=\psi(x,t)\cdot\chi(\abs{\omega}/M)\) and proves
(\texttt{sec:moser-s33}) that each of the three second-derivative-of-magnitude
terms --- \(\nabla^2\log m\) from the drift, \(\Delta m\) inside
\(\Dt\chi\), and the \(\chi'\)-cutoff terms --- integrates by parts onto
first-order factors bounded by \(4\abs{\nabla m}/M\), and is then absorbed by
Young's inequality against \(\nu\abs{\nabla^2\ehat}^2\). The level-set boundary
contributes nothing, because \(\chi\) vanishes below \(M/4\).

The session also opens with a non-circularity remark that is worth naming,
because it is a direct answer to the audit. The pincer is a contradiction
argument at a putative \emph{first} singularity, so every estimate runs on the
smooth interval \([0,T)\); Defect~D4 --- the Leray--Hopf circularity --- cannot
recur on this route by construction.

The output is \texttt{prop:first-rung} (line~5740), a Caccioppoli inequality
for the angular energy, \PROVEDMOD{four enumerated routine items}. It is here
that the logarithm formally enters the machinery: \texttt{lem:local-enstrophy-typeI}
(line~5701) states the local enstrophy-gradient bound
\(\nu\iint_{Q(2\rho)}\abs{\nabla\omega}^2\lesssim M^2\rho^3\Lfac\), and defines
\(\Lfac=\log(e+\norm{u}_{H^3}/M)\) at line~5712.\footnote{Line numbers in this
chapter and in Chapters~\ref{ch:unification}--\ref{ch:non-implication} are
checked against the proof document as it stands after \Sn{48}. They differ
from the numbers quoted in Chapters~\ref{ch:inheritance} and \ref{ch:walls}
for every object after line~6900, because \Sn{48} inserted an amendment
paragraph into the \Sn{39} ledger and a repair into
\texttt{prop:profile-extraction}: \texttt{prop:transfer-audit} moved from
6941 to 6955, \texttt{rem:s39-distance} from 7670 to 7703,
\texttt{rem:s41-shortfall} from 8395 to 8428, and
\texttt{prop:s44-linfty-route} from 9209 to 9242. The labels are unchanged and
are what should be cited.}

\subsection{\texorpdfstring{\Sn{34}: an a priori bound, and the first sighting of the wall}{S34: an a priori bound, and the first sighting of the wall}}

\Sn{34} opens with a corrigendum to \Sn{33} --- a dimensional re-audit of the
previous session's own section, performed \emph{before} building on it, which
found four defects and fixed them. Two inhomogeneous terms had the wrong power
of \(\rstar\); the corrected terms vanish as \(\rstar\to0\), so the fix is
favourable. A third item, previously filed as routine bookkeeping, was
reclassified: the \(\nu w\ehat\) term of \eqref{eq:pincer-dir} yields the
critical quadratic \(\nu\iint w^2\), which is the standard critical
nonlinearity and needs the smallness hypothesis \(\sstar\le\varepsilon_0\).
That is where H1's own hypothesis enters its own proof, as it should.

The session's positive result is the first unconditional quantitative output
of the whole pincer programme.

\begin{statement}[A priori coherence bound --- \texttt{thm:apriori-sigma}, line~5873]
\label{st:apriori}
\PROVED{} Near a Type-I singularity,
\[
  \langle\sstar\rangle
  \;=\; (\rstar)^{-3}\iint_{Q(\rstar/2)} w\,\mathbf 1_{\{\abs{\omega}\ge M/2\}}
  \;\le\; C(C_I)\,\frac{\Lfac}{\nu}.
\]
\end{statement}

The proof tests \eqref{eq:pincer-mag} against \(m\eta^2\), which puts the
good-sign sink \(\nu\iint m^2w\eta^2\) on the left. The key device is the
pointwise consequence of the split identity,
\(m^2\abs{\nabla\ehat}^2=\abs{\nabla\omega}^2-\abs{\nabla m}^2\le\abs{\nabla\omega}^2\),
which sends every transition-annulus term straight to the local
enstrophy-gradient budget with \emph{no absorption argument at all}. That is
\texttt{lem:grad-omega-split}, the audit's survivor, doing load-bearing work
fifteen sessions later.

The coherence deficit near a Type-I singularity is therefore a priori bounded,
at order \(\log\). What H1 needs is \(\sstar\le\varepsilon_0\). The gap between
``automatic'' and ``needed'' is the explicit factor \(C\Lfac/(\nu\varepsilon_0)\)
--- which converts the programme's smallness requirement into a \emph{decay}
question for the \(\sstar\) dynamics. That reformulation is what the rest of
the programme is about.

The same session states the negative half sharply.
\texttt{prop:crude-K} (line~5903) proves that the crude route to absorbing the
coupling term \(K\) re-enters at top order multiplied by \(\Lfac\): it ``misses
closing the loop by exactly one logarithm''. And it names the hypothesis that
would close it, \texttt{conj:K-refined} (line~5917): the correlation constant
\[
  c_K \;=\; \frac{\langle\abs{\nabla m}^2 w\rangle_E}
                 {\langle\abs{\nabla m}^2\rangle_E\,\langle w\rangle_E},
  \qquad E=Q(\rstar)\cap\{m\ge M/4\},
\]
satisfies \(c_K\le C/\Lfac\) near a Type-I singularity. The mechanism is
equipartition: by the split identity, \(\abs{\nabla m}^2\) and \(m^2w\) are
competitors for the same \(\abs{\nabla\omega}^2\) budget, so only strong
\emph{positive} correlation could defeat absorption.

That conjecture is the origin of four sessions of numerics and of
Chapter~\ref{ch:numerical-lessons}'s hardest lesson. Its measurement history
is in Chapter~\ref{ch:layer4}; what it could and could not decide is Wall~6.

\subsection{\texorpdfstring{\Sn{35}: the second rung, and a term that did not exist}{S35: the second rung, and a term that did not exist}}

\texttt{prop:second-rung} (line~5986) disposes of the critical quadratic
\(\nu\iint w^2\eta^2\) by the pointwise split \(w^2=f^{10/3}w^{1/3}\) with
\(f=\abs{\nabla\ehat}\eta\), the parabolic embedding
\(L^\infty L^2\cap L^2H^1\hookrightarrow L^{10/3}(Q)\), and Young. Under
\(\sstar\le\varepsilon_0\) the inhomogeneous remainder is
\(O((\Lfac\varepsilon_0)^{5/2}\rstar)\), superlinear in \(\varepsilon_0\) and
therefore harmless. The entire nonlinearity of the direction equation is
subordinate to \(\sstar\)-smallness.

The same session recorded \texttt{prop:K-self} (line~6026), explicitly marked
\emph{provisional}, claiming \emph{two} good terms on the left of a weighted
identity. \Sn{36} found that the second cancels exactly
(\texttt{lem:w2-cancel}, line~6166) and it was withdrawn. That episode is
Wall~7 and is discussed in Chapter~\ref{ch:unification}; it is mentioned here
because \Sn{35} is also the session in which the programme wrote down, in the
proof document itself, an honest answer to the owner's direct question of
whether any of this proves existence. The answer recorded was \emph{no}:
weak-solution global existence is Leray 1934, the open question is smoothness,
and the near-term realistic target is the conditional disorder-barrier
theorem, ``a real contribution, not the Prize''.

\section{\texorpdfstring{\Sn{38}: Regime A is impossible, and what that costs}{S38: Regime A is impossible, and what that costs}}
\label{sec:regime-a}

By the end of \Sn{37} the programme's open analytical content had been reduced
to two items (Chapter~\ref{ch:unification}): the annulus residual, then
believed mechanical, and one deep gap --- forcing \(\sstar(t)\to0\). \Sn{37}
had also left one apparent escape: Regime~A, in which the global enstrophy
\(\Omega=\norm{\omega}_{L^2}^2\) stays bounded near the putative singularity,
and in which \texttt{prop:regime-A} (line~6430) gives depletion with
\(\theta=2/5<1/2\) unconditionally, using only the a priori bound and no decay
at all.

\Sn{38} checked that escape route first, as a fact-check, before deriving
anything --- and closed it.

\begin{statement}[\texttt{prop:regime-A-impossible}, line~6514]
\label{st:regime-a-impossible}
\IMPOSSIBLE{} At any genuine blow-up time \(T\),
\(\limsup_{t\to T^-}\Omega(t)=\infty\). Hence Regime~A cannot occur.
\end{statement}

The argument is short and standard, which is why it is decisive. By
Fujita--Kato local well-posedness, a solution with bounded \(H^1\) norm extends
for a time depending only on that norm, so blow-up forces
\(\limsup\norm{u}_{H^1}=\infty\); the energy inequality makes
\(\norm{u}_{L^2}\) non-increasing, so \(\norm{\nabla u}_{L^2}\to\infty\); and
Biot--Savart on \(\T^3\) turns that into \(\Omega\to\infty\).

\begin{obsv}[A proof that narrows its own programme]
\label{obs:regime-a-cost}
\Sn{38} is the clearest instance in this book of a result that is
simultaneously a genuine advance and a loss of ground. Before it, the
programme had two ways to reach H2: prove \(\sstar\)-decay, or establish
bounded enstrophy. After it, there is one. The document states the consequence
without softening: ``\(\sstar\)-decay is not one of several paths to H2 ---
given H1, it is \textbf{necessary}''.

\smallskip
The same proposition is what makes Wall~6 provable. A resolved numerical
simulation is a bounded-enstrophy object throughout, so the regime the
programme's conjectures speak about is not merely one the runs failed to
visit; it is not of the type a resolved run represents. That inference is
recorded in Chapter~\ref{ch:layer4} and again in
Chapter~\ref{ch:numerical-lessons}, and it is the reason four sessions of
\(c_K\) measurement were never going to decide anything.
\end{obsv}

\Sn{38} also declined to force a result. The natural next step --- assembling
the proved first rung, second rung and amended \(K\)-absorption into a genuine
scale recursion for \(\sstar(\rho)\) as \(\rho\to0\) --- was attempted and left
\OPEN{}, with the reason stated: whether the effective coefficient is
subcritical (\(<1\), self-improving regardless of the starting value) or
critical (\(\ge1\), needing a priori smallness, which is the fate that
befell \S19 and \S20) was not determined. A preliminary scaling check on
leading terms showed consistent powers of \(M\), which the document notes is
necessary but not sufficient. The session log gives the reason for the
restraint plainly: the programme had by then made three mid-session
corrections on exactly this class of multi-term scaling argument, and a fourth
was not wanted.

\section{The arc, session by session}
\label{sec:pincer-ledger}

\begin{table}[ht]
\centering
\footnotesize
\begin{tabular}{@{}l p{0.40\textwidth} p{0.35\textwidth}@{}}
\toprule
Session & What it added & Status of the addition today \\
\midrule
\Sn{32} & Direction equation re-derived; \(\sstar\) defined and proved
  scale-invariant; \(\sstar\le16\sigma\); marginality; H1, H2 and the pincer
  stated
  & \PROVED{} except H1 and H2, both \OPEN{} \\[3pt]
\Sn{33} & Level-set Moser cutoff; three integrations by parts;
  \(\Lfac\) enters at \texttt{lem:local-enstrophy-typeI}; first rung
  & \PROVEDMOD{four routine items}, all discharged in \Sn{36} \\[3pt]
\Sn{34} & Corrigendum to \Sn{33} (four defects, fixed favourably);
  \(\langle\sstar\rangle\le C\Lfac/\nu\); crude \(K\) fails by one
  logarithm; \texttt{conj:K-refined}
  & \PROVED{}; \PROVED{} sharp; \OPEN{} \\[3pt]
\Sn{35} & Second rung (critical quadratic subordinated to \(\varepsilon_0\));
  \(K\)-self-absorption identity, marked provisional; honest
  distance-to-Prize remark
  & \PROVED{}; second good term \WITHDRAWN{} at \Sn{36} \\[3pt]
\Sn{36} & C1, C2 verified; Gram negativity; \(w^2\) cancellation;
  routine items discharged; annulus residual isolated
  & Chapter~\ref{ch:unification} \\[3pt]
\Sn{37} & Localized Constantin--Fefferman derived; Regime A/B dichotomy;
  H2 and \(\sstar\)-decay unified
  & Chapter~\ref{ch:unification} \\[3pt]
\Sn{38} & Regime~A \IMPOSSIBLE{}; \(\sstar\)-decay necessary; closure target
  posed and deliberately left open
  & \IMPOSSIBLE{}; the target is \OPEN{} \\
\bottomrule
\end{tabular}
\caption[The pincer arc]{Sessions \Sn{32}--\Sn{38}. Three of the seven
sessions begin by correcting the previous one. Two of the seven produce a
proof that closes off one of the programme's own options.}
\label{tab:pincer-arc}
\end{table}

\begin{obsv}[What the pincer is, and what it is not]
\label{obs:pincer-ceiling}
Even fully closed --- H1 proved, H2 proved --- the pincer yields
Statement~\ref{st:pincer}: \emph{every Type-I singularity maintains persistent
geometric disorder}. That is not regularity. It excludes coherent Type-I
singularities and says nothing whatever about disordered ones, and nothing at
all about Type-II. Turning it into ``no Type-I singularity'' requires the
further step of forcing \(\sstar(t)\to0\), which is
\texttt{target:disorder-depletion} and is \OPEN{}; turning \emph{that} into
the Prize requires Type-II exclusion, which \texttt{rem:s39-distance} places
outside these methods entirely. Chapter~\ref{ch:profile-rigidity} states the
ceiling in the document's own words, and Chapter~\ref{ch:where-it-stands}
keeps it in view.
\end{obsv}
